\documentclass[../main.tex]{subfiles}
\begin{document}

\section{Probleme}

\subsection{Problema Festival (ONI 2025, clasele 11-12)}
\label{problem:festival}

\href{https://kilonova.ro/problems/3720}{enunț}
$\bullet$
\hyperref[code:festival]{sursă}

Începem cu precizarea că există și soluții mai eficiente la această problemă (\href{https://kilonova.ro/submissions/1161191}{exemplu}), bazate pe tehnica \href{https://robert1003.github.io/2020/01/31/cdq-divide-and-conquer.html}{Divide et Impera CDQ}. Dar iată și soluția de manual cu structuri 2D offline.

\subsubsection*{Formula de recurență}

Este relativ evident că festivalurile formează un DAG. Nodurile corespund festivalurilor, iar între un festival $A$ și un festival $B$ există o muchie dacă și numai dacă Andrei poate ajunge de la $A$ la $B$. Fiind un DAG, graful are o sortare topologică, de exemplu dacă sortăm festivalurile crescător după $t$. De aici rezultă și o formulă implementabilă în $\bigoh(n^2)$. Fie $D_i$ satisfacția maximă pentru a vizita un șir de festivaluri care se încheie la festivalul $i$. Atunci:

$$D_i = S_i + \max_{j} \{ D_j \ | \ \text{Andrei poate ajunge de la $j$ la $i$} \}$$

Putem figura grafic un festival $(x, t)$ ca pe un punct $(x, y=t)$. Atunci zona din plan din care Andrei poate ajunge la un festival are forma din figura \ref{fig:offline-2d-festival}:

\import{./figures}{offline-2d-festival.tex}

Astfel reducem problema la una de căutări și actualizări într-o structură de date. Pentru fiecare festival $i$ în ordine crescătoare a lui $y$, calculăm valoarea festivalului ca $S_i$ plus valoarea maximă din zona de acoperire. Apoi notăm această valoare la coordonatele festivalului.

\subsubsection*{Reducerea la structuri de date 2D}

Operațiile pe structuri de date diagonale pot fi complicate. De aceea, să încercăm să reducem zona de acoperire la dreptunghiuri.

Dacă împărțim zona mergînd pe linia punctată care coboară din $P$, obținem două trapeze. Să analizăm trapezul din dreapta. Nu putem extinde pur și simplu acest trapez la o zonă dreptunghiulară, deoarece am include și punctul $D$, de la care Andrei nu poate ajunge la $P$. În schimb, să procedăm astfel. Vom înlocui fiecare punct $(x, y)$ cu punctul $(x, y + x)$. Cu alte cuvinte, pe fiecare coloană $i$ translatăm punctele în sus cu $i$ pași. Astfel, zona trapezoidală devine dreptunghică. Desigur, vom introduce în zona de acoperire și niște puncte cu $y$ negativ, dar aceasta nu denaturează rezultatele, întrucît nu există festivaluri cu timp negativ.

Întrebarea „care este satisfacția maximă în trapezul drept al lui $P(x, y)$?” devine o interogare de maxim în dreptunghiul $(x, -\infty) - (x + d, y + x)$.

Pentru trapezul stîng, dorim ca punctele din stînga să urce mai mult decît cele din dreapta, deci vom inversa translația, înlocuind fiecare punct $(x, y)$ cu punctul $(x, y - x)$. Remarcăm că aceasta poate duce la coordonate $y$ negative, dar vom vedea odată ce definitivăm structura de date că valorile negative nu ne deranjează. Vom folosi două structuri de date, cîte una pentru fiecare tip de translație.

Întrebarea „care este satisfacția maximă în trapezul stîng al lui $P(x, y)$?” devine o interogare de maxim în dreptunghiul $(x - d, -\infty) - (x, y - x)$.

Desigur, $D_i$ este maximul valorilor din cele două trapeze. După calculare, trebuie să notăm valoarea $D_i$ obținută la $(x, y + x)$ în structura dreaptă și la $(x, y - x)$ în structura stîngă.

\subsubsection*{Structura necesară}

Acum știm ce operații trebuie să admită structura dorită:

\begin{enumerate}
  \item Notarea unei valori într-un punct.
  \item Găsirea valorii maxime dintr-un dreptunghi care se extinde în jos pînă la $-\infty$.
\end{enumerate}

Coordonatele respectă $0 \leq x \leq 10^9$ și $-2 \cdot 10^9 \leq y \leq 2 \cdot 10^9$. Ca să putem folosi structuri 2D offline, trebuie să normalizăm valorile $x$. Este nevoie de atenție la interogări, căci valorile $x - d$ și $x + d$ nu înseamnă „$\pm d$ poziții în structură”, ci înseamnă „cea mai mică poziție pe care se găsește o valoare cel puțin egală cu $x - d$” și respectiv „cea mai mare poziție pe care se găsește o valoare cel mult egală cu $x + d$.

Coordonatele $y$ nu trebuie normalizate.

Deoarece avem nevoie de maxim pe interval, vom folosi un arbore de intervale ca structură primară. În structurile secundare avem nevoie de maxim pe prefix, în condițiile în care valorile pot doar să crească, deci putem folosi și AIB-uri. Ca detaliu de implementare, putem folosi AIB-uri indexate de la 0 sau, dacă dorim să le folosim pe cele mai naturale, indexate de la 1, trebuie să nu uităm să adăugăm cîte o santinelă cu $y = -\infty$ în fiecare AIB.

De aici încolo, implementarea este tipică.

\subsection{Problema Medwalk, revizitată (Lot 2025)}
\label{problem:medwalk-fenwick}

\href{https://kilonova.ro/problems/3790}{enunț}
$\bullet$
\hyperref[code:medwalk-fenwick]{sursă}

Am discutat această problemă și în \hyperref[problem:medwalk]{capitolul de arbori de intervale}. Am găsit o soluție bazată pe un AINT de seturi, construit peste valorile din matrice. Să vedem acum una de 5 ori mai rapidă, probabil cea pe care a dorit-o comisia.

Primele observații se mențin. Decuplăm (conceptual) matricea în doi vectori, unul cu minimele și unul cu maximele de pe fiecare coloană. Pentru a minimiza medianul (scorul unui interval $[l, r]$), căutăm un drum minim lexicografic. Acesta va consta din minimele de pe intervalul $[l, r]$ și din minimul maximelor. Medianul acestei mulțimi va fi una din trei valori posibile (logica deciziei este simplă):

\begin{itemize}
  \item fie minimul maximelor;
  \item fie medianul minimelor;
  \item fie elementul anterior medianului minimelor.
\end{itemize}

Pentru cazul (1), minimul maximelor îl menținem într-un aint simplu, cu actualizări punctuale și cu interogări de minim pe interval. Pentru cazurile (2) și (3) trebuie să răspundem la interogări de al $k$-lea element pe interval. Aici introducem următoarea soluție, constînd dintr-un \href{https://usaco.guide/plat/2DRQ?lang=cpp#offline-2d-bit}{AIB 2D offline} construit peste minimele din matrice.

Dorim să căutăm binar al $k$-lea element. Bunăoară, prima dată ne întrebăm: este el mai mare decît $V_{max} / 2$? Dacă da, îl căutăm între $V_{max} / 2$ și $V_{max}$. Dacă nu, îl căutăm între 1 și $V_{max} / 2$. În general, pentru plaja de valori curentă, $[x, y]$, vom calcula mijlocul $m = (x + y) / 2$ și îi vom adresa structurii de date întrebarea: cîte valori între $x$ și $m$ apar la intrare pe poziții între $l$ și $r$? Dacă răspunsul este $\geq k$, atunci al $k$-lea element are valoarea între $x$ și $m$. Altfel, el are valoarea între $m$ și $y$.

Să încercăm o structură naivă: o matrice binară uriașă $A$ de dimensiuni $n \times V_{max}$, unde $A_{p,v}$ reține 1 dacă valoarea $v$ apare în matrice la poziția $p$, 0 altfel. Atunci răspunsul la întrebarea „cîte valori între $v_1$ și $v_2$ apar pe poziții între $l$ și $r$?” este suma dreptunghiului $[l, r] \times [v_1, v_2]$. Sună bine, dar memoria nu ne permite această matrice.

De aceea, vom comprima fiecare coloană într-un AIB pornind de la următoarea observație. Coloana $v$ va reține 1, oricînd pe durata întregului program, doar pe liniile $p$ cu proprietatea că poziția $p$ în vectorul de minime va avea, cel puțin după o operație, valoarea $v$. De aceea, în primă fază simulăm operațiile și colectăm toate minimele posibile pe toate pozițiile, care vor fi cel mult $n + q$. (De aici provine partea de \textit{offline} din nume.) Pe fiecare coloană $v$ stocăm:

\begin{itemize}
  \item O listă a pozițiilor posibile pentru valori 1, ordonate crescător.
  \item Un AIB de 0/1 construit peste această listă.
\end{itemize}

De exemplu, dacă valoarea $9$ este măcar o dată minimă pe pozițiile 102, 109, 110, 113, atunci pe coloana 9 stocăm vectorul $[102, 109, 110, 113]$ și un AIB de 0/1 cu 4 poziții disponibile.

Dar facem mai mult decît atît, căci v-am promis un AIB 2D. \emoji{slightly-smiling-face} Vectorul de coloane este el însuși un AIB. Coloana $v$ nu reține doar pozițiile aparițiilor lui $v$, ci ale oricărei valori din $(v - p, v]$, unde $p$ este cea mai mare putere a lui 2 mai mică sau egală cu $v$. În particular, cele 4 valori 102, 109, 110, 113 exemplificate mai sus nu apar doar în lista coloanei 9, ci și în listele coloanelor $10, 12, 16, 32, 64, \dots$.

\import{./figures}{offline-2d-medwalk.tex}

Să observăm că fiecare poziție apare în lista unui număr logaritmic de coloane. De aceea, suma lungimilor tuturor listelor (și a tuturor AIB-urilor) este $\bigoh(n \log V_{max})$.

Pe această structură putem răspunde în timp $\bigoh(n \log n \log V_{max})$ la interogări de al $k$-lea element.

\end{document}
